\scanpage{6}
Considérons alors
\[
 fg\xrightarrow{\alpha}\mathrm{id} \tag*{(2.8)}
\]
ce qui fournit un homomorphisme d'adjonction
\[
 \begin{aligned}
 &\inserted{v\longmapsto\alpha\deleted{\unsure{(Y)}}\circ f(v):}\\
 &\deleted{\unsure{\mathrm{Hom}(fX,Y)\longleftarrow}}\alpha':\mathrm{Hom}_A(X,g(Y))\longrightarrow\mathrm{Hom}_B(f(X),Y)
 \end{aligned} \tag*{(2.9)}
\]
\deleted{\unsure{par}} On va définir aussi
\[
 \deleted{v\longmapsto\alpha\circ f(\unsure{u|v})} \tag*{\deleted{(2.9')}}
\]
\[
 \deleted{\unsure{g}}\beta:\mathrm{id}_K\longrightarrow gf \tag*{(2.\deleted{\unsure{9}}10)}
\]
\unsure{par|en} \unsure{prenant|posant}, pour tt $Y=(X,u)$ dans $K$,
\diagram{\begin{tikzcd}[ampersand replacement=\&]
 X\arrow[r,"u"]\arrow[d,"u"'] \& \varphi(X)\arrow[d,dashed,"\varphi(u)"] \& \unsure{\leftarrow}Y,\\
 \varphi(X)\arrow[r,"\lambda(X)"'] \& \varphi^2(X) \& g\deleted{\unsure{g}}f(Y)
\end{tikzcd}}
la commutativité étant assurée par la \uuline{condition} \textcircled{\scriptsize (1.\deleted{\unsure{9}}10)} sur les objets de $K$. On a ainsi un hom. d'adjonction en l'autre sens
\[
 \beta':\mathrm{Hom}_B(fX,Y)\longrightarrow\mathrm{Hom}_A(X,g(Y))
 \qquad v\longmapsto g(v)\circ\beta(X) \tag*{(2.11)}
\]
On vérifie que ces hom. sont \uline{inverses l'un de l'autre}. En effet, $\beta'\alpha'=\mathrm{id}$ résulte facilement de \textcircled{\scriptsize 1.5}, et $\alpha'\beta=\mathrm{id}$ \inserted{résulte de \textcircled{\scriptsize (1.\deleted{\unsure{10}}9)}.}
\deleted{équivaut à ceci : Pour \unsure{tout} objet $(X,u)$ de $A$, le composé}
\diagram{\begin{tikzcd}[ampersand replacement=\&,every arrow/.append style={dash pattern=on 2pt off 0pt}]
 \deleted{X}\arrow[r,"{\deleted{u}}"] \& \deleted{\varphi(X)}\arrow[r,"{\deleted{\alpha(X)}}"] \& \deleted{X}
\end{tikzcd}}
\deleted{est \uline{l'identité}. Or \unsure{on a} (1.9) $u=\mathrm{id}_{\varphi(X)}u
 \overset{1.5}{=}\varphi(\alpha(X))\lambda(X)u\overset{1.9}{=}\varphi(\alpha(X))\varphi(u)u
 =\varphi(\alpha(X)u)$ et (1.5) $\varphi(u)\circ u=\lambda(X)\circ u:X\longrightarrow\varphi^2(X)$
 d'où, en composant à gauche avec $\lambda(X)$
 $\underbrace{\lambda(X)\varphi(u)}\circ u=\underbrace{\lambda(X)\lambda(X)}u$}

Donc $f$ et $g$ sont adjoints l'un de l'autre, et \uline{redonnent bien sûr} \deleted{\unsure{par le}} $(\varphi,\alpha,\lambda)$ \uline{par le procédé du} N° 1.

Le foncteur $f$ est conservatif et satisfait \unsure{?|c|e} \ldots
% archive sheet number: 5
